\documentclass[../Main.tex]{subfiles}

\begin{document}
\section{Review of Relativity}
\subsection{Lorentz Transformation}
In the special theory of relativity we use spacetime coordinates:
\begin{equation*}
    X = \begin{pmatrix}ct \\ x \\ y \\ z\end{pmatrix} = \begin{pmatrix}ct \\ \vec{x}\end{pmatrix}
\end{equation*}
and reference the coordinates of $X$ as $X^\mu$. We use Greek letters for the components of 4-vectors. These run from $0$ to $3$. Roman indices will run from $1$ to $3$ to only reference spatial dimensions.

$X$ is the position 4-vector. Under a Lorentz transformation (LT) from frame $S$ to $S'$, it transforms according to
\begin{equation}
    X' = \Lambda X
    \label{eqnLorentzTransform}
\end{equation}
In index notation, we would find $X^{\prime \mu} = \indx{\Lambda}{\mu}{\nu} X^\nu$.

$\Lambda$ is any matrix in the Lorentz group. We are particularly interested in the Lorentz boost, where $S'$ moves relative to $S$ with speed $v < c$ in the $x$ direction.
\begin{equation*}
    \Lambda =
    \begin{pmatrix}
        \gamma & -\gamma \beta & 0 & 0 \\
        -\gamma\beta & \gamma & 0 & 0 \\
        0 & 0 & 1 & 0 \\
        0 & 0 & 0 & 1
    \end{pmatrix}
\end{equation*}
where $\beta = \frac{v}{c} \in (-1, 1)$, and:
\begin{equation*}
    \gamma = \frac{1}{\sqrt{1 - \beta^2}}
\end{equation*}
is the \underline{Lorentz factor}.
\subsection{Four-Vectors}
Four-vectors will be denoted with captial Roman letters. Any 4-vector $V$ transforms according to \eqnref{eqnLorentzTransform}, $V' = \Lambda V$. We use the convention that any such vector can be written as:
\begin{equation*}
    V = \begin{pmatrix}V^0 \\ \vec{v}\end{pmatrix}
\end{equation*}
The pseudo-inner product of 4-vectors is:
\begin{equation*}
    V \cdot W = \vec{v} \cdot \vec{w} - V^0 W^0 = \eta_{\mu \nu} V^\mu W^\nu
\end{equation*}
where $\eta$ is the \underline{Minkowski metric}, and it is an isotropic tensor. It has components:
\begin{equation*}
    \eta_{\mu\nu} = \begin{pmatrix}
        -1 & 0 & 0 & 0 \\
        0 & 1 & 0 & 0 \\
        0 & 0 & 1 & 0 \\
        0 & 0 & 0 & 1
    \end{pmatrix}
    = \operatorname{diag}(-1, 1, 1, 1)
\end{equation*}
in any intertial frame.

Then the inner product is invariant under a Lorentz transformation. We see this:
\begin{align*}
    V' \cdot W' &= \eta_{\mu\nu}' V^{\prime \mu} W^{\prime \nu} \\
    &= \eta_{\mu \nu} \indx{\Lambda}{\mu}{\rho} V^\rho = \indx{\Lambda}{\nu}{\rho} W^\sigma \\
    &= \eta_{\rho \sigma} V^\rho W^\sigma \\
    &= V \cdot W
\end{align*}
because the LT matrix satisfies:
\begin{equation*}
    \eta_{\mu \nu} \indx{\Lambda}{\mu}{\rho} \indx{\Lambda}{\nu}{\sigma} = \eta_{\rho \sigma}
\end{equation*}
or, in matrix notation,
\begin{equation*}
    \Lambda^T \eta \Lambda = \eta
\end{equation*}
which implies $\det(\Lambda) = \pm 1$.

We can define the square of a $4$-vector to be $V \cdot V$:
\begin{equation*}
    V \cdot V = -V^0 V^0 + |\vec{v}|^2
\end{equation*}
and we say:
\begin{itemize}
    \item $V$ is \underline{timelike} if $V \cdot V < 0$,
    \item $V$ is \underline{spacelike} if $V \cdot V > 0$,
    \item $V$ is \underline{null} or \underline{lightlike} if $V \cdot V = 0$.
\end{itemize}
The invariant spacetime interval between two events (an event is a point and a time):
\begin{equation*}
    ds^2 = dX \cdot dX = -c^2 dt^2 + dx^2 + dy^2 + dz^2
\end{equation*}

\begin{definition}{4-velocity}
    The \underline{4-velocity} of a massive particle is:
    \begin{equation*}
        U = \frac{dX}{d\tau}
    \end{equation*}
\end{definition}
where $\tau$ is the proper time (time experienced by the particle).

Since $X = (ct, \vec{x})$,
\begin{align*}
    U &= \frac{dt}{d\tau} \begin{pmatrix}c \\ \vec{u}\end{pmatrix} \\
    &= \gamma_{\vec{u}} \begin{pmatrix}c \\ \vec{u}\end{pmatrix}    
\end{align*}
where $\vec{u}$ is the normal 3-velocity and $\gamma_{\vec{u}}$ is the Lorentz factor for the motion:
\begin{equation*}
    \gamma_{\vec{u}} = \frac{1}{1 - \frac{|\vec{x}|^2}{c^2}}
\end{equation*}
As for all 4-vectors, the dot product is invariant:
\begin{align*}
    U \cdot U &= \left(\frac{dt}{d\tau}\right)^2 (-c^2 + |\vec{u}|^2) \\
    &= -c^2
\end{align*}
Therefore this is invariant and timelike.

\begin{definition}{4-momentum}
    The \underline{4-momentum} of a massive particle with rest mass $m$ is:
    \begin{equation*}
        P = mU = m\gamma_{\vec{u}} \begin{pmatrix}c \\ \vec{u}\end{pmatrix} = \begin{pmatrix}\frac{E}{c} \\ \vec{p}\end{pmatrix}
    \end{equation*}
    where $\vec{p}$ is the 3-momentum.
\end{definition}
We know that the inner product is invariant:
\begin{align*}
    P \cdot P &= -\left(\frac{E}{c}\right)^2 + |\vec{p}|^2 \\
    &= -m^2 c^2
\end{align*}
which is consistent with $E^2 = m^2 c^4 + |\vec{p}|^2 c^2$.

In the Newtonian limit $v \ll c$,
\begin{align*}
    E &= m \gamma_{\vec{u}} c^2 \\
    &\approx m\left(1 + \frac12 \frac{|\vec{u}|^2}{c^2}\right) c^2 \\
    &= mc^2 + \frac12 m |\vec{u}|^2
\end{align*}
and in the rest frame where $\vec{u} = \vec0$, we simply have the rest energy $E = mc^2$.

Massless particles such as photons travel at the speed of light (in any inertial frame). They have $E = |\vec{p}|c$ and their 4-velocity is null.
\section{Covectors and Tensors}
We understand a 4-vector to be some $X$ in the relevant vector space, with components $X^\mu$ (superscript indices).
\begin{definition}{Covector}
    A \underline{covector} is an element $Y$ of the dual space, with components $Y_\mu$ (subscript indices).
\end{definition}
Define:
\begin{equation*}
    Y_\mu = \frac{\partial f}{\partial X^\mu}
\end{equation*}
and this is the derivative of a scalar field $f(X)$ in spacetime. This is a covector, and so we want to know how this transforms under a Lorentz transformation. By the chain rule,
\begin{align}
    \frac{\partial f}{\partial X^{\prime \mu}} &= \frac{\partial X^\nu}{\partial X^{\prime \mu}} \frac{\partial f}{\partial X^\nu} \nonumber \\
    Y_\mu' &= \indx{(\Lambda^{-1})}{\nu}{\mu} Y_\nu \label{eqnLorentzTransformDerivative}
\end{align}
where $(\Lambda^{-1})$ is the inverse Lorentz transformation matrix. \Eqnref{eqnLorentzTransformDerivative} is the general transformation law for covectors. Written without the indices,
\begin{equation*}
    Y' = \Lambda^{-1} Y
\end{equation*}
\begin{definition}{Tensor}
    A \underline{(4-)tensor} $T$ of type $(r, s)$ is an object whose $4^{r+s}$ elements transform as:
    \begin{equation}
        \indx{T}{\prime \mu_1 \cdots \mu_r}{\nu_1 \cdots \nu_s} = \indx{\Lambda}{\mu_1}{\rho_1} \cdots \indx{\Lambda}{\mu_r}{\rho_r} \indx{(\Lambda^{-1})}{\sigma_1}{\nu_1} \cdots \indx{(\Lambda^{-1})}{\sigma_s}{\nu_s} \indx{T}{\rho_1 \cdots \rho_r}{\sigma_1\cdots\sigma_s}
        \label{eqnLorentzTransformTensor}
    \end{equation}
    That is, a LT matrix for each superscript index and an inverse LT matrix for each subscript index.
\end{definition}
\begin{remarks}
    \item Any ordering of the up and down indices is possible. For example, a tensor of type $(1, 1)$, $A_\mu\hphantom{}^\nu$
    \item A $(1, 0)$ tensor is a 4-vector.
    \item A $(0, 1)$ tensor is a 4-covector.
    \item A $(0, 0)$ tensor is a scalar (Lorentz invariant).
\end{remarks}
\begin{propositions}{
        
        \label{propsTensorProps}
    }
    \item A linear combination of $(r, s)$ tensors is also an $(r, s)$ tensor. \label{propTensorLinearCombo}
    \item The \underline{outer product} of tensors of type $(r_1, s_1)$ and $(r_1, s_2)$ gives a tensor of type $(r_1 + r_2, s_1 + s_2)$. \label{propTensorOuterProd}
    \item The \underline{derivative} of a tensor field of type $(r, s)$ is another tensor field of type $(r, s+1)$.
        \begin{equation*}
            A_\mu\indx{}{\nu \rho}{\sigma} = \partial_\mu \indx{B}{\nu \rho}{\sigma}
        \end{equation*}
        where $\partial_\mu = \frac{\partial}{\partial X^\mu}$. \label{propTensorDerivative}
    \item The \underline{contraction} of an $(r, s)$ tensor on a pair of its indices produces an $(r-1, s-1)$ tensor. This only works when an up index is contracted with a down index. \label{propTensorContract}
\end{propositions}
The summation convention on spacetime requires any repeated index to appear once up and once down. This is because a contraction is an application of an element of the dual space to an element of the vector space. For example, the invalid expression $X^\mu X^\mu$ is not Lorentz invariant and so is not valid as a scalar.

\begin{definition}{Isotropic}
    A tensor is \underline{isotropic} if its components are the same in every inertial frame.
\end{definition}
\begin{remark}
    An example of this is the Minkowski metric $\eta_{\mu \nu}$. This is an isotropic tensor of type $(0, 2)$. The inverse of the Minkowski metric is $\eta^{\mu\nu}$ which is an isotropic tensor of type $(2, 0)$. Both have the same components, but they have different roles.
\end{remark}
The Kronecker Delta is an isotropic tensor of type $(1, 1)$. $\delta = diag(1,1, 1, 1)$.

The Levi-Civita epsilon symbol is defined by:
\begin{equation*}
    \epsilon_{\mu\nu\rho\sigma} =
    \begin{cases}
        1 & \text{if $(\mu, \nu, \rho, \sigma)$ is an even permutation of $(0,1,2,3)$} \\
        -1 & \text{if $(\mu, \nu, \rho, \sigma)$ is an odd permutation of $(0,1,2,3)$} \\
        0 & \text{otherwise}
    \end{cases}
\end{equation*}
and it is an isotropic pseudo-tensor of type $(0, 4)$. A pseudo-tensor satisfies the same transformation law as a tensor but with an extra factor of $\det(\Lambda)$.
\begin{definition}{Tensor equation}
    A \underline{tensor equation} is an equation in which every term is a tensor of the same type (transforms identically under a LT). Such an equation is said to be \underline{covariant} and is valid in any inertial frame.
\end{definition}
\end{document}